% Options for packages loaded elsewhere
\PassOptionsToPackage{unicode}{hyperref}
\PassOptionsToPackage{hyphens}{url}
\documentclass[
  ignorenonframetext,
]{beamer}
\newif\ifbibliography
\usepackage{pgfpages}
\setbeamertemplate{caption}[numbered]
\setbeamertemplate{caption label separator}{: }
\setbeamercolor{caption name}{fg=normal text.fg}
\beamertemplatenavigationsymbolsempty
% remove section numbering
\setbeamertemplate{part page}{
  \centering
  \begin{beamercolorbox}[sep=16pt,center]{part title}
    \usebeamerfont{part title}\insertpart\par
  \end{beamercolorbox}
}
\setbeamertemplate{section page}{
  \centering
  \begin{beamercolorbox}[sep=12pt,center]{section title}
    \usebeamerfont{section title}\insertsection\par
  \end{beamercolorbox}
}
\setbeamertemplate{subsection page}{
  \centering
  \begin{beamercolorbox}[sep=8pt,center]{subsection title}
    \usebeamerfont{subsection title}\insertsubsection\par
  \end{beamercolorbox}
}
% Prevent slide breaks in the middle of a paragraph
\widowpenalties 1 10000
\raggedbottom
\AtBeginPart{
  \frame{\partpage}
}
\AtBeginSection{
  \ifbibliography
  \else
    \frame{\sectionpage}
  \fi
}
\AtBeginSubsection{
  \frame{\subsectionpage}
}
\usepackage{iftex}
\ifPDFTeX
  \usepackage[T1]{fontenc}
  \usepackage[utf8]{inputenc}
  \usepackage{textcomp} % provide euro and other symbols
\else % if luatex or xetex
  \usepackage{unicode-math} % this also loads fontspec
  \defaultfontfeatures{Scale=MatchLowercase}
  \defaultfontfeatures[\rmfamily]{Ligatures=TeX,Scale=1}
\fi
\usepackage{lmodern}
\ifPDFTeX\else
  % xetex/luatex font selection
\fi
% Use upquote if available, for straight quotes in verbatim environments
\IfFileExists{upquote.sty}{\usepackage{upquote}}{}
\IfFileExists{microtype.sty}{% use microtype if available
  \usepackage[]{microtype}
  \UseMicrotypeSet[protrusion]{basicmath} % disable protrusion for tt fonts
}{}
\makeatletter
\@ifundefined{KOMAClassName}{% if non-KOMA class
  \IfFileExists{parskip.sty}{%
    \usepackage{parskip}
  }{% else
    \setlength{\parindent}{0pt}
    \setlength{\parskip}{6pt plus 2pt minus 1pt}}
}{% if KOMA class
  \KOMAoptions{parskip=half}}
\makeatother
\setlength{\emergencystretch}{3em} % prevent overfull lines
\providecommand{\tightlist}{%
  \setlength{\itemsep}{0pt}\setlength{\parskip}{0pt}}
\usepackage{bookmark}
\IfFileExists{xurl.sty}{\usepackage{xurl}}{} % add URL line breaks if available
\urlstyle{same}
\hypersetup{
  hidelinks,
  pdfcreator={LaTeX via pandoc}}

\author{\texorpdfstring{}{}}
\date{}

\begin{document}

\begin{frame}{The Alchemical Forges: Los vs.~The Royal Mint}
\protect\phantomsection\label{the-alchemical-forges-los-vs.-the-royal-mint}
If Urizen is the bureaucratic Master of the Mint---the Newtonian
administrator calculating ratios, imposing geometric limits, and
standardizing the infinite human soul into weighed coins---then Los is
the divine blacksmith, the alchemist of the spirit. The opposition
between these two figures encodes one of the most structurally coherent
economic critiques in the entire Romantic canon.

\begin{block}{The Mint on Tower Hill: Engines of Homogenization}
\protect\phantomsection\label{the-mint-on-tower-hill-engines-of-homogenization}
In \emph{Jerusalem}, Los is constantly depicted laboring at his roaring
furnaces. Blake utilizes explicitly metallurgical and alchemical imagery
to describe this work. However, whereas the Royal Mint on Tower
Hill---relocated in 1810 to a purpose-built facility on Little Tower
Hill designed by James Johnson and Robert Smirke---was an industrial
fortress of steam-powered homogenization, Los's forges are designed to
produce living, differentiated, irreducibly unique human forms (Eden).
The Mint's industrial transformation was itself a symptom of the
Urizenic revolution: Matthew Boulton's steam-powered coining presses,
installed at the Soho Mint in Birmingham and later adopted at the Royal
Mint, could strike coins with mechanical precision at high speed,
eliminating the human hand from the minting process entirely and
replacing the artisan's judgment with calibrated, replicable force.

The Mint melts down the various, chaotic silver and gold of the world
and stamps it with the sovereign's head---a singular image of authority.
It is a process of extraction and homogenization that parallels exactly
the Urizenic cosmogony described in \emph{The First Book of Urizen}
(1794), where the tyrant ``form'd golden compasses / And began to
explore the Abyss'' {[}@blake1795urizen{]}. Los's furnaces, conversely,
operate on a principle of spiritual alchemy: taking the fragmented,
degraded, atomic material of the fallen world and fusing it back
together through the heat of prophetic imagination into the ``Human Form
Divine.'' This opposition encodes a fundamental first-principles insight
into the nature of value itself: the Mint's operation is
\emph{subtractive} (reducing diverse human labor to a single metallic
denominator), whereas Los's forge is \emph{additive} (synthesizing
fragmented experience into integrated wholes).

The metallic imagery in \emph{Jerusalem} is pervasive and deliberate. On
Plate 12, Blake describes the spiritual architecture of
Golgonooza---Los's prophetic city, the anti-London---in explicitly
metallurgical terms:

\begin{quote}
``The stones are pity, and the bricks, well wrought affections, Enameld
with love \& kindness, \& the tiles engraven gold, Labour of merciful
hands: the beams \& rafters are forgiveness: The mortar \& cement of the
work, tears of honesty: the nails, And the screws \& iron braces, are
well wrought Blandishments'' (\emph{Jerusalem}, Plate 12, lines 30--34;
Erdman p.~155) {[}@erdman1988complete{]}
\end{quote}

Here the ``tiles engraven gold'' and ``iron braces'' serve a function
opposed to the Mint's gold: they are construction materials for a city
built from human affect---pity, kindness, forgiveness, honesty---rather
than the weight of bullion. Further, on the same plate, Blake names
``the golden Looms of Cathedron'' and ``the golden mills of Los'' (Plate
12, lines 46--47; Erdman p.~156), establishing that Los possesses his
own golden instruments, but ones devoted to weaving and milling the
fabric of human reality rather than stamping coins of abstract
equivalence. The adjective ``golden'' in Blake is thus bifurcated:
Urizen's gold is dead and hoarded; Los's gold is living, laboring, and
regenerative.
\end{block}

\begin{block}{The Boulton Press and the Artist's Body}
\protect\phantomsection\label{the-boulton-press-and-the-artists-body}
Blake explicitly contrasts the mechanical, steam-powered presses of
early industrial capitalism---embodied most visibly by the Boulton steam
presses---with the vital, bodily labor of the artist. The production of
true value requires sweat, breath, and the integrated labor of the
entire human body, not the alienated pull of a lever. As Thompson
documents, Blake's hostility to industrial machinery was not Luddite
nostalgia but a precise philosophical objection: the machine severs the
connection between the laborer's creative will and the product of their
hands, reproducing the structure of Urizenic alienation at the level of
daily lived experience {[}@thompson1993witness{]}.

\begin{quote}
``And Los's Furnaces howl loud; living: self-moving: lamenting With Fury
\& Despair, \& they dictate the relations of States'' (\emph{Jerusalem},
Plate 73) {[}@erdman1988complete{]}
\end{quote}

The furnaces are not inert mechanisms; they are \emph{living},
\emph{self-moving}, \emph{lamenting}---they participate in the suffering
of the human condition even as they labor to transcend it. This is
Blake's radical proposition: that the sites of genuine
production---artistic, spiritual, political---are necessarily sites of
anguish and fury, not the automated efficiency of the Mint.
\end{block}

\begin{block}{The Anti-Coin: Blake's Illuminated Printing as Economic
Subversion}
\protect\phantomsection\label{the-anti-coin-blakes-illuminated-printing-as-economic-subversion}
For Blake, the true economy is not dictated by the Bank of England or
the bullion brokers of the Royal Exchange, but by the ``Furnaces of
Los.'' The artist's labor---specifically Blake's invention of
illuminated printing via relief-etching (c.~1788)---becomes the only
valid site of socio-metaphysical reproduction. Blake's relief-etching
process is the economic inverse of the Mint's stamping dies: where
Newton's Urizenic 1717 Great Recoinage \emph{stamped} an image of
sovereign authority onto a blank disc to impose fungible value, Blake
\emph{dissolved} the surface of the copper plate with acid to reveal a
unique, irreproducible spiritual topography. The Mint standardizes the
infinite human soul into fixed bimetallic weights; Blake's acid bath
frees the living form from the dead copper. Each illuminated plate is an
anti-coin: singular, non-fungible, resistant to the metric logic of
exchange {[}@poovey2008genres{]}. Its value does not inhere in the
weight of its copper but in the density of the visionary experience it
encodes.

The choice facing humanity is clear: be minted as a passive coin under
Urizen, or forge one's own infinite value in the roaring fires of Los
{[}@blake1804jerusalem{]}. As Erdman argues, this is not merely artistic
preference but political economy: Blake's mode of production constitutes
a complete alternative to the industrial system, one in which the
producer retains full sovereignty over the means, process, and meaning
of their labor {[}@erdman1954prophet{]}.
\end{block}
\end{frame}

\end{document}
